




The Petz recovery map~\cite{petz1986sufficient, petz1988sufficiency} is a powerful tool with broad applications in quantum science~\cite{barnum2002reversing, ng2010simpleapproach, beigi2016decoding, chen2020entanglement, lauten2022approx, nakayama2023petz}.
In particular, it is known that the map demonstrates nearly optimal performance when it reverses the effect of noise~\cite{barnum2002reversing, beigi2016decoding}.
The Petz recovery map $\cR_{\sigma, \cF}^{\sB \rarr \sA}$ is defined for a reference state $\sigma^\sA$ and a quantum channel $\cF^{\sA\rarr\sB}$ as $\cR_{\sigma, \cF}^{\sB \rarr \sA} = \cQ_{\sigma^{1/2}}^\sA\circ(\cF^{\sA\rarr\sB})^\dag\circ\cQ_{\cF(\sigma)^{-1/2}}^\sB$, where $(\cF^{\sA\rarr\sB})^\dag$ is an adjoint map of $\cF^{\sA\rarr\sB}$ and $\cQ_{\mu}^\sC = \mu^\sC(\cdot)\mu^\sC$ for an operator $\mu^\sC$.
The adjoint map $(\cF^{\sA\rarr\sB})^\dag$ of $\cF^{\sA\rarr\sB}$ is defined as the map satisfying $\tr[L^\sB\cF^{\sA\rarr\sB}(M^\sA)]=\tr[(\cF^{\sA\rarr\sB})^\dag(L^\sB)M^\sA]$ for any operators $L^\sB$ and $M^\sA$.


Quantum algorithms for the Petz recovery map have been studied in Refs.~\cite{gilyen2022petzmap, biswas2023noiseadapted}.
These works represent significant progress toward implementation of the map, but their approaches require, in addition to the use of $\sigma^\sA$, access to a Stinespring unitary of the channel $\cF^{\sA\rarr\sB}$ and its inverse, which is challenging in practice.


We here introduce a method to realize the Petz recovery map using only $\sigma^{\sA}$ and $\cF^{\sA\rarr\sB}$; that is, access to a Stinespring unitary is not required.
The implementation of $(\sigma^{\sA})^{1/2}$ and $\big(\cF^{\sA \rarr \sB}(\sigma^\sA)\big)^{-1/2}$ via block-encodings, directly using the density matrix exponentiation and the QSVT with multiple uses of $\sigma^{\sA}$ and $\cF^{\sA \rarr \sB}$, is straightforward (see, e.g., Refs.~\cite{gilyen2022petzmap, gilyen2022improvedfidelity, wang2023SampletoQuary, wang2024TimeEfficientEntEstSamp} for references).
Hence, we focus our investigation only on the non-trivial part $(\cF^{\sA\rarr\sB})^\dag$.
To this end, we make use of the following proposition, in which $(\cP_{\ket{0}}^{\bC\rarr\sG})^\dag$ denotes a linear map (not a quantum channel), defined as $(\cP_{\ket{0}}^{\bC\rarr\sG})^\dag(\cdot) = \bra{0}^\sG \cdot \ket{0}^\sG$.





\begin{proposition}[Implementation of a Stinespring unitary via the Uhlmann transformation algorithm]
\label{prop:stinespring unitary via uhlmann}
Let $\delta \in (0, 1)$. For any quantum channel $\cF^{\sA\rarr\sB}$, there exists a quantum algorithm that implements a quantum channel $\cG^\sS$ satisfying 
\begin{align}
    \f{1}{2}\big\|\cG^\sS \circ \cP_{\ket{0}}^{\bC\rarr\sG} - \cU_\cF^\sS \circ \cP_{\ket{0}}^{\bC\rarr\sG}\big\|_\diamond \leq \delta, 
\end{align}
where $\cU_\cF^\sS$ is a Stinespring unitary of $\cF^{\sA\rarr\sB}$ on $\sS = \sA\sG = \sB\sE$, with the dimension of $\sE$ taken as $d_\sE = d_\sA d_\sB$.
The algorithm uses the channel $\cF^{\sA \rarr \sB}$ a total of $\nu$ times, where $\nu = \til{\cO}\Big(\f{d_\sA^4 d_\sB}{\delta^2}\min\Big\{\f{1}{\tau_{\rm min}^2}, \f{d_\sA^6d_\sB^2}{\delta^4}\Big\}\Big)$, and $\tau_{\rm min}$ is the minimum non-zero eigenvalue of the Choi--Jamio\l kowski state $\tau^{\sR\sB} = \cF^{\sA\rarr\sB}(\Phi^{\sR\sA})$.

Moreover, there exists a quantum algorithm that implements a quantum channel $\cG_{\rm inv}^\sS$ satisfying 
\begin{align}
    \f{1}{2}\big\|(\cP_{\ket{0}}^{\bC\rarr\sG})^\dag\circ\cG_{\rm inv}^\sS - (\cP_{\ket{0}}^{\bC\rarr\sG})^\dag\circ(\cU_\cF^\sS)^\dag\big\|_\diamond \leq \delta,
\end{align}
also using $\cF^{\sA\rarr\sB}$ a total of $\nu$ times.

In both algorithms, $\cO\big(\log{(d_\sA d_\sB)}\big)$ qubits suffice at any one time.

\end{proposition}


Proposition~\ref{prop:stinespring unitary via uhlmann} will be proven at the end of this section.





For an algorithmic implementation of the Petz recovery map, we follow the protocol in Ref.~\cite{gilyen2022petzmap}.
The Petz recovery map can be rephrased as
\begin{align}
\label{eq:petz exp form rephrase}
    \cR_{\sigma, \cF}^{\sB \rarr \sA}
    &= d_\sE\cQ_{\sigma^{1/2}}^\sA\circ(\cP_{\ket{0}}^{\bC\rarr\sG})^\dag \circ(\cU_\cF^{\sS})^\dag \circ\cP_\pi^{\bC\rarr\sE}\circ\cQ_{\cF(\sigma)^{-1/2}}^\sB,
\end{align}
where $\sS = \sA\sG = \sB\sE$.
Since the size of the system $\sE$ is arbitrarily chosen, we set $\sE = \sR'\sB'$ with dimension $d_\sE = d_\sA d_\sB$. Note that $d_\sR = d_\sA$.
While post-selection by $\ket{0}^\sG$ may seem necessary, it can in fact be removed using the robust oblivious amplitude amplification~\cite{gilyen2019qsvt, Gilyn2019thesis} by iteratively applying $U_\cF^\sS$ and $(U_\cF^\sS)^\dag$ a total of $v = \cO\big(\sqrt{d_\sA d_\sB/\lambda_{\rm min}}\big)$ times~\cite{gilyen2022petzmap}.
Here, $\lambda_{\rm min}$ denotes the minimum non-zero eigenvalue of $\cF^{\sA\rarr\sB}(\sigma^\sA)$.
Hence, an approximate Petz recovery map can be deterministically implemented.




In the $v$ repetitions of $U_\cF^\sS$ and $(U_\cF^\sS)^\dag$ in the oblivious amplitude amplification, when we approximate these unitaries by $\cG^\sS$ and $\cG_{\rm inv}^\sS$ within an error $\delta$, respectively, the total error accumulates to $v\delta$.
By rescaling $\delta$ as $\delta = \epsilon/v$, we implement the approximation of the Petz recovery map with error at most $\epsilon$ in the diamond norm, using $\cF^{\sA \rarr \sB}$ the following total number of times:
\begin{align}
\label{eq:sampcomp Petz}
    v\til{\cO}\Big(\f{d_\sA^4 d_\sB}{(\epsilon/v)^2 }\min\Big\{\f{1}{\tau_{\rm min}^2}, \f{d_\sA^6d_\sB^2}{(\epsilon/v)^4}\Big\}\Big) =\til{\cO}\Big(\f{d_\sA^{11/2} d_\sB^{5/2}}{\epsilon^2  \lambda_{\rm min}^{3/2}}\min\Big\{\f{1}{\tau_{\rm min}^2}, \f{d_\sA^8 d_\sB^4}{\epsilon^4\lambda_{\rm min}^2}\Big\}\Big).
\end{align}


To be precise, implementing $\cQ_{\sigma^{1/2}}^\sA$ and $\cQ_{\cF(\sigma)^{-1/2}}^\sB$ additionally requires access to $\sigma^\sA$ and $\cF^{\sA\rarr\sB}$.
However, as mentioned, since these can be straightforwardly realized using the density matrix exponentiation and the QSVT, we do not provide a detailed discussion here.







Other approaches to implementing the Petz recovery map $\cR_{\sigma, \cF}^{\sB\rarr\sA}$ using only the channel $\cF^{\sA\rarr\sB}$ and the state $\sigma^\sA$ may be possible. 
In fact, since the map is independent of the choice of the system $\sE$ (see Eq.~\eqref{eq:petz exp form rephrase}), the recent results in Ref.~\cite{tang2025conjugatequerieshelp} based on the random purification techniques appear to be applicable. 
This approach would achieve computational costs of roughly the square of those in Ref.~\cite{gilyen2022petzmap} divided by the accuracy, although it would require collective measurements over an exponential number of Choi--Jamio\l kowski states of $\cF^{\sA\rarr\sB}$.
We leave a detailed investigation of this approach for future work.







We now prove \Cref{prop:stinespring unitary via uhlmann}.


\begin{proof}[Proof of \Cref{prop:stinespring unitary via uhlmann}]
    


We consider a purified state of the Choi--Jamio\l kowski state $\tau^{\sR\sB} = \cF^{\sA\rarr\sB}(\Phi^{\sR\sA})$, which is given by $\ket{\tau}^{\sR\sB\sE} = U_\cF^{\sS}\ket{\Phi}^{\sR\sA}\ket{0}^{\sG}$, where $U_\cF^{\sS}$ is a Stinespring unitary of $\cF^{\sA\rarr\sB}$ and $\sS = \sA\sG = \sB\sE$. Without loss of generality, we take $d_\sR = d_\sA$. 
Since $\tau^{\sR} = \pi^{\sR}$, we observe that $U_\cF^\sS$ corresponds to the Uhlmann unitary transforming $\ket{\Phi}^{\sR\sA}\ket{0}^{\sG}$ to $\ket{\tau}^{\sR\sB\sE}$.
Thus, using our Uhlmann transformation algorithm for $\ket{\Phi}^{\sR\sA}\ket{0}^{\sG}$ and $\ket{\tau}^{\sR\sB\sE}$, the unitary $U_\cF^{\sS}$ can be implemented.


To generate a purified state of $\tau^{\sR\sB}$, we use the canonical purification: $\ket{\tau_{\rm c}}^{\sR\sB\sR'\sB'} = (\sqrt{\tau^{\sR\sB}} \otimes \bI^{\sR'\sB'})\ket{\Gamma}^{\sR\sB\sR'\sB'}$. Note that the environment $\sE$ is chosen to be $\sR'\sB'$ in the canonical purification. 
Thus, $\sS = \sA\sG = \sB\sR'\sB'$.
From \Cref{prop:canonical purification}, to obtain a state $\til{\tau}_{\rm c}^{\sR\sS} = \mathtt{CanonicalPurification}(\tau^{\sR\sB}; \delta_1)$ that approximates $\ket{\tau_{\rm c}}^{\sR\sS}$ with the error $\delta_1$ in the trace distance: $\f{1}{2}\big\|\til{\tau}_{\rm c}^{\sR\sS} - \ketbra{\tau_{\rm c}}{\tau_{\rm c}}^{\sR\sS}\big\|_1 \leq \delta_1$, we use $f$ samples of $\tau^{\sR\sB}$, where $f = \til{\cO}\Big(\f{d_\sA d_\sB}{\delta_1}\min\Big\{ \f{1}{\tau_{\rm min}^2}, \f{d_\sA^2 d_\sB^2}{\delta_1^4}\Big\}\Big)$, and $\tau_{\rm min}$ is the non- zero minimum eigenvalue of $\tau^{\sR\sB}$.



Next, we apply the Uhlmann transformation algorithm with $\ket{\sigma}^{\sR\sS} = \ket{\Phi}^{\sR\sA}\ket{0}^{\sG}$ and $\til{\tau}_{\rm c}^{\sR\sS}$ to obtain an approximation of the unitary $U_\cF^\sS$.
At this point, a careful analysis is needed, since the Uhlmann transformation algorithm realizes a partial isometry $W_\cF^\sS$, not the full unitary $U_\cF^\sS$.
However, this does not cause any critical issue, as we will show below.



The relation between the partial isometry $W_\cF^\sS$ and the Stinespring unitary $U_\cF^\sS$ is that $W_\cF^\sS$ is embedded in the part of $U_\cF^\sS$ specified by projections $\ketbra{0}{0}^\sG\otimes\bI^\sA$ and $\Pi_\cF^{\sB\sR'\sB'} = V_\cF^{\sA\rarr\sB\sR'\sB'}(V_\cF^{\sA\rarr\sB\sR'\sB'})^\dag$, where $V_\cF^{\sA\rarr\sB\sR'\sB'}$ is a Stinespring isometry of $\cF$. That is, 
\begin{align}
        U_\cF^\sS
        = \hspace{1mm}
\begin{blockarray}{ccc}
 & \ketbra{0}{0}^\sG\otimes\bI^\sA &  & \vspace{1mm}\\
\begin{block}{c(cc)}
  \Pi_\cF^{\sB\sR'\sB'} \hspace*{1mm} & W_\cF^\sS & \hspace{0mm} * \hspace{3mm}\\
   \hspace*{1mm} & * & \hspace{0mm} * \hspace{3mm}\\
\end{block}
\end{blockarray}\hspace{2.5mm}.
\end{align}
Then, the following relation holds:
\begin{align}
    U_\cF^\sS\ket{0}^\sG 
    &= V_\cF^{\sA\rarr\sB\sR'\sB'} \\
    &= V_\cF^{\sA\rarr\sB\sR'\sB'}(V_\cF^{\sA\rarr\sB\sR'\sB'})^\dag V_\cF^{\sA\rarr\sB\sR'\sB'} \\
    &= \Pi_\cF^{\sB\sR'\sB'} U_\cF^{\sS} (\ketbra{0}{0}^\sG\otimes\bI^\sA) \ket{0}^\sG \\
    & = W_\cF^\sS \ket{0}^\sG.
\end{align}
Hence, when we fix an input on the system $\sG$ as $\ket{0}^\sG$, the actions of them are identical.


Let $U_{\rm ideal}^{\sS\sH\hat{\sR}\hat{\sS}}$ be an exact block-encoding unitary of $\ketbra{\sigma}{\tau_{\rm c}}^{\hat{\sR}\hat{\sS}} \otimes W_\cF^\sS$.
Then, a quantum channel $\cJ^{\sS\sH\hat{\sR}\hat{\sS}} = \mathtt{UhlmannPurifiedSample}(\ket{\sigma}^{\sR\sS}, \til{\tau}_{\rm c}^{\sR\sS}; \delta_2, \diamond)$ satisfies
\begin{align}
\label{inteq:122}
    \f{1}{2}\big\|\cJ^{\sS\sH\hat{\sR}\hat{\sS}}\circ\cP_{\ket{0}}^{\bC\rarr\sH}\circ\cP_{\til{\tau}_{\rm c}}^{\bC\rarr\hat{\sR}\hat{\sS}} - \cU_{\rm ideal}^{\sS\sH\hat{\sR}\hat{\sS}} \circ \cP_{\ket{0}\ket{\tau_{\rm c}}}^{\bC\rarr\sH\hat{\sR}\hat{\sS}}\big\|_\diamond \leq \delta_2 + (2u+1)\delta_1,
\end{align}
where $u=\cO\big(d_\sA\log{(1/\delta_2)}\big)$.
The number of samples of $\til{\tau}_{\rm c}^{\sR\sB\sR'\sB'}$ and $\ket{\sigma}^{\sR\sS}$ is given by 
\begin{align}
    w =\cO\big(d_\sA^2\big(\log{(1/\delta_2)}\big)^2/\delta_2\big).
\end{align}
We here used the fact that the minimum non-zero singular value of $\sqrt{\pi^{\sR}}\sqrt{\tau_{\rm c}^{\sR}} = \pi^\sR$ is $1/d_\sA$.
The case where one input to $\mathtt{UhlmannPurifiedSample}$ is a mixed state $\til{\tau}_{\rm c}^{\sR\sS}$ can be handled by a similar argument to that in \Cref{sec:proof eval sample mixed}.

Note that $U_{\rm ideal}^{\sS\sH\hat{\sR}\hat{\sS}}\ket{0}^{\sH}\ket{\tau_{\rm c}}^{\hat{\sR}\hat{\sS}}\ket{0}^\sG = \ket{0}^\sH\ket{\sigma}^{\hat{\sR}\hat{\sS}}W_\cF^\sS\ket{0}^\sG = \ket{0}^\sH\ket{\sigma}^{\hat{\sR}\hat{\sS}}U_\cF^\sS\ket{0}^\sG$.
When we define a quantum channel $\cG^\sS$ as $\cG^\sS = \tr_{\sH\hat{\sR}\hat{\sS}} \circ \cJ^{\sS\sH\hat{\sR}\hat{\sS}} \circ \cP_{\ket{0}}^{\bC\rarr\sH}\circ\cP_{\til{\tau}_{\rm c}}^{\bC\rarr\hat{\sR}\hat{\sS}}$, we have that
\begin{align}
    &\f{1}{2}\big\|\cG^\sS\circ\cP_{\ket{0}}^{\bC\rarr\sG} - \cU_\cF^\sS\circ\cP_{\ket{0}}^{\bC\rarr\sG}\big\|_\diamond \notag \\
    &\leq \f{1}{2}\big\|\cJ^{\sS\sH\hat{\sR}\hat{\sS}}\circ\cP_{\ket{0}}^{\bC\rarr\sH}\circ\cP_{\til{\tau}_{\rm c}}^{\bC\rarr\hat{\sR}\hat{\sS}}\circ\cP_{\ket{0}}^{\bC\rarr\sG} - \cU_{\rm ideal}^{\sS\sH\hat{\sR}\hat{\sS}}\circ\cP_{\ket{0}}^{\bC\rarr\sH}\circ\cP_{\ket{\tau_{\rm c}}}^{\bC\rarr\hat{\sR}\hat{\sS}}\circ\cP_{\ket{0}}^{\bC\rarr\sG}\big\|_\diamond \\
    &\leq \f{1}{2}\big\|\cJ^{\sS\sH\hat{\sR}\hat{\sS}}\circ\cP_{\ket{0}}^{\bC\rarr\sH}\circ\cP_{\til{\tau}_{\rm c}}^{\bC\rarr\hat{\sR}\hat{\sS}} - \cU_{\rm ideal}^{\sS\sH\hat{\sR}\hat{\sS}} \circ \cP_{\ket{0}\ket{\tau_{\rm c}}}^{\bC\rarr\sH\hat{\sR}\hat{\sS}}\big\|_\diamond \\
    &\leq \delta_2 + (2u+1)\delta_1,
\end{align}
where we used Eq.~\eqref{inteq:122} in the third inequality.
Rescaling $\delta_1$ and $\delta_2$ as $\delta_1 = \delta/\big(2(2u+1)\big)$ and $\delta_2 = \delta/2$, the over all error is bounded from above by $\delta$.

In this algorithm, the number of uses of the quantum channel $\cF^{\sA\rarr\sB}$ is evaluated as
\begin{align}
    \nu 
    &= fw \\
    \label{inteq:19}
    &= \til{\cO}\Big(\f{d_\sA^4 d_\sB}{\delta^2}\min\Big\{\f{1}{\tau_{\rm min}^2}, \f{d_\sA^6d_\sB^2}{\delta^4}\Big\}\Big).
\end{align}


Constructing a quantum algorithm for approximating $(U_\cF^\sS)^\dag$ can be achieved straightforwardly using the fact that the Uhlmann unitary from $\ket{\tau_{\rm c}}^{\sR\sS}$ to $\ket{\sigma}^{\sR\sS} = \ket{\Phi}^{\sR\sA}\ket{0}^\sG$ coincides with $(U_\cF^\sS)^\dag$.
Hence, a quantum channel $\cJ_{\rm inv}^{\sS\sH\hat{\sR}\hat{\sS}} = \mathtt{UhlmannPurifiedSample}(\til{\tau}_{\rm c}^{\sR\sS}, \ket{\sigma}^{\sR\sS}; \delta_2, \diamond)$ satisfies that
\begin{align}
\label{inteq:123}
    &\f{1}{2}\big\|\cJ_{\rm inv}^{\sS\sH\hat{\sR}\hat{\sS}}\circ\cP_{\ket{0}\ket{\sigma}}^{\bC\rarr\sH\hat{\sR}\hat{\sS}} - (\cU_{\rm ideal}^{\sS\sH\hat{\sR}\hat{\sS}})^\dag \circ \cP_{\ket{0}\ket{\sigma}}^{\bC\rarr\sH\hat{\sR}\hat{\sS}}\big\|_\diamond 
    \leq \delta_2 + 2u\delta_1.
\end{align}
Note that, compared to the previous case for $U_\cF^\sS$, the first and second arguments of $\mathtt{UhlmannPurifiedSample}$ are swapped.


We observe that  $\bra{0}^\sG(U_{\rm ideal}^{\sS\sH\hat{\sR}\hat{\sS}})^\dag\ket{0}^{\sH}\ket{\sigma}^{\hat{\sR}\hat{\sS}} = \ket{0}^\sH\ket{\tau_{\rm c}}^{\hat{\sR}\hat{\sS}}\bra{0}^\sG(W_\cF^\sS)^\dag = \ket{0}^\sH\ket{\tau_{\rm c}}^{\hat{\sR}\hat{\sS}}\bra{0}^\sG (U_\cF^\sS)^\dag$.
Defining a quantum channel $\cG_{\rm inv}^\sS$ as $\cG_{\rm inv}^\sS = \tr_{\sH\hat{\sR}\hat{\sS}} \circ \cJ_{\rm inv}^{\sS\sH\hat{\sR}\hat{\sS}} \circ \cP_{\ket{0}\ket{\sigma}}^{\bC\rarr\sH\hat{\sR}\hat{\sS}}$, we see that it satisfies
\begin{align}
    &\f{1}{2}\big\|(\cP_{\ket{0}}^{\cC\rarr\sG})^\dag\circ\cG_{\rm inv}^\sS
    - (\cP_{\ket{0}}^{\cC\rarr\sG})^\dag\circ(\cU_\cF^\sS)^\dag\big\|_\diamond \\
    &\leq \f{1}{2}\big\|(\cP_{\ket{0}}^{\cC\rarr\sG})^\dag\circ \cJ_{\rm inv}^{\sS\sH\hat{\sR}\hat{\sS}}\circ\cP_{\ket{0}\ket{\sigma}}^{\bC\rarr\sH\hat{\sR}\hat{\sS}} - (\cP_{\ket{0}}^{\cC\rarr\sG})^\dag\circ(\cU_{\rm ideal}^{\sS\sH\hat{\sR}\hat{\sS}})^\dag\circ\cP_{\ket{0}\ket{\sigma}}^{\bC\rarr\sH\hat{\sR}\hat{\sS}}\big\|_\diamond \\
    &\leq \f{1}{2}\big\|(\cP_{\ket{0}}^{\cC\rarr\sG})^\dag\|_\diamond \big\|\cJ_{\rm inv}^{\sS\sH\hat{\sR}\hat{\sS}}\circ\cP_{\ket{0}\ket{\sigma}}^{\bC\rarr\sH\hat{\sR}\hat{\sS}} - (\cU_{\rm ideal}^{\sS\sH\hat{\sR}\hat{\sS}})^\dag\circ\cP_{\ket{0}\ket{\sigma}}^{\bC\rarr\sH\hat{\sR}\hat{\sS}}\big\|_\diamond \\
    &\leq \delta_2 + 2u\delta_1,
\end{align}
where in the second inequality, we used the fact that, for any linear map $\cL$ and $\cM$, $\|\cL\circ\cM\|_\diamond \leq \|\cL\|_\diamond\|\cM\|_\diamond$ holds~\cite{watrous2018TheoryQI}, and we used Eq.~\eqref{inteq:123} and $\|(\cP_{\ket{0}}^{\bC\rarr\sG})^\dag\|_\diamond = 1$ in the third inequality. 
Rescaling $\delta_1$ and $\delta_2$ as $\delta_1 = \delta/4u$ and $\delta_2 = \delta/2$, the over all error is bounded from above by $\delta$. 
The number of uses of $\cF^{\sA\rarr\sB}$ is given in the same order as Eq.~\eqref{inteq:19}.

In both cases of approximating $U_\cF$ and $(U_\cF)^\dag$, the procedures are conducted sequentially, and thus, $\cO\big(\log{(d_\sA d_\sB)}\big)$ qubits suffice at any one time.


\end{proof}


